\subsection{Integral higher periods from the vertical condition alone}
\label{arch136:vertical-periods}

\paragraph{Sheet and logical scope.}
This gives an application of
ARCH-136-ASSOCIATIVE-\allowbreak COMPARISON.
It uses only the vertical hypothesis for the common-target route.
It concerns the vertical and horizontal bounds on cote~136,
physical pp.~117,119, footers 116,118, and the common-source
question on physical p.~121, footer 120, supplied scans
\texttt{ARCHIVE\_136\_scan117.tex},
\texttt{ARCHIVE\_136\_scan119.tex},
\texttt{ARCHIVE\_136\_scan121.tex}.
The proof uses the already constructed associative bar resolution,
not a new strict PD comparison. No input A, M, or outside
preprint is used.

\paragraph{ARCH-136-VERTICAL-PERIODS.}
Let \(A\) be a simplicial nonnegative weight-graded differential
\(S=\Gamma_{\mathbf Z}(T)\)-algebra. Suppose that, for each \(d\),
there is a bound \(V_A(d)\) such that at all weights
\(w\geq V_A(d)\), uniformly in the simplex,
\[
 S^w\xrightarrow{\sim}H^0(A_n^{w,*}),\qquad
                  H^q(A_n^{w,*})=0\quad(1\leq q\leq d).
 \tag{136.VP1}
\]
No horizontal acyclicity is assumed. There are canonical
homomorphisms, natural in \(A\) and in simplicial pairs \((X,Y)\),
\[
 \operatorname{per}_A:
 H^j(A^{w,*}(X,Y))
       \longrightarrow H^j(X,Y;S^w)
 \qquad (0\leq j\leq d,\ w\geq V_A(d+1)).
 \tag{136.VP2}
\]
Here \(A(X,Y)\) denotes the kernel of restriction to \(Y\).
These maps preserve products whenever the input and output
bidegrees lie in the stated comparison ranges.

If \(A_{\mathbf Q}\) is horizontally and vertically admissible,
then, after increasing the bound in each fixed degree,
the canonical rational comparison sends every integral
class in \(H^j(A(X,Y))\) to the image of
\[
                 H^j(X,Y;S^w)\longrightarrow
                    H^j(X,Y;S_{\mathbf Q}^w).
 \tag{136.VP3}
\]
In particular, all ordinary periods on finite simplicial spheres,
or on a simplex relative to its boundary, are integral in a
uniform weight tail in every fixed degree.

\emph{Proof.}
Use the representable bar resolution \(P_*(X)\) with its
Alexander--Whitney diagonal from
ARCH-136-ASSOCIATIVE-COMPARISON. For a pair, put
\[
                  P_*(X,Y)=P_*(X)/P_*(Y).
 \]
Each bar term of \(P_*(X)\) is a direct sum of modules
\(\mathbf Z[\Delta[n]]\), and the inclusion for \(Y\subseteq X\)
is the inclusion of a subfamily of those summands. This remains
true at every iteration of the free representable functor.
Hence every \(P_r(X,Y)\) is projective in simplicial abelian
groups. The quotient is a projective resolution of
\(\mathbf Z[X]/\mathbf Z[Y]\): exactness follows by quotienting
the two augmented bar resolutions.

Define the same convolution complex as in the inherited
construction,
\[
 K_A(X,Y)^{w,n}
 =\bigoplus_{r+q=n}
       \operatorname{Hom}_{s\mathrm{Ab}}
                (P_r(X,Y),A^{w,q}).
 \tag{136.VP4}
\]
The Alexander--Whitney diagonal descends to the relative
quotient. It gives the usual nonunital relative product and,
with the inherited totalization signs, natural associative
differential algebra maps
\[
 A(X,Y)\longrightarrow K_A(X,Y)
                         \longleftarrow K_{\underline S}(X,Y).
 \tag{136.VP5}
\]
The right-hand complex computes ordinary relative cohomology
with the coefficient \(S^w\).

The right-hand arrow in (136.VP5) is an isomorphism on
cohomology through degree \(d\) under (136.VP1).
Indeed projectivity of \(P_r(X,Y)\) identifies the vertical
cohomology of its Hom complex with
\(\operatorname{Hom}(P_r(X,Y),H^q(A^{w,*}))\).
The coefficient row is therefore the only nonzero row in
the required finite range. The first-quadrant spectral sequence
has finite diagonals, and (136.VP1) through \(d+1\) is a safe
bound for all the asserted degrees. This argument uses no
horizontal condition on \(A\) and no flatness of its terms.
Define (136.VP2) by the left-hand map followed by the inverse
of this right-hand cohomology isomorphism.

Naturality follows from the functorial bar resolution and
the uniqueness of the inverse. The convolution maps are
multiplicative; therefore their induced maps and inverse
cohomology identifications preserve products whenever the
three relevant ranges hold. For relative products with
different subspaces, the same argument uses the usual pair
diagonal into the corresponding relative tensor product.

There is a natural diagram of (136.VP5) under
\(A\to A_{\mathbf Q}\) and \(S\to S_{\mathbf Q}\).
When \(A_{\mathbf Q}\) is admissible, the rational left-hand
arrow is also a cohomology isomorphism in a uniform range
by the full associative comparison. For pairs this follows
either from the same projective resolution or from absolute
comparison and the pair exact sequences: horizontal
acyclicity makes restriction surjective in the required
form degrees by the normalized free relative-chain extension.
The natural diagram gives (136.VP3).

For infinite simplicial sets, no assertion that tensoring
with \(\mathbf Q\) commutes with the products in a cochain
complex is needed. The diagram uses the actual natural
coefficient maps between the integral and rational Hom
complexes. Thus the argument applies to arbitrary pairs,
not merely finite ones.
\hfill\(\square\)

\paragraph{ARCH-136-INITIAL-TENSOR-PERIODS.}
Let \(C\in\mathcal A_{\mathbf Z}\), and let \(B\) be the
standard integral theory. Then
\[
                            A=B\otimes_S C
 \]
has all the ordinary integral periods in (136.VP3),
uniformly in every fixed degree, without any flatness
hypothesis on \(C\) and without a strict map \(C\to B\).
The same holds after the incoming FIRST-SKELETON replacement
of \(C\).

\emph{Proof.}
In each simplex the ordered-coordinate PD contraction of
\(B_n\) onto \(S\) is \(S\)-linear and weight-preserving.
Tensoring this actual homotopy with \(C_n\) identifies
the vertical cohomology of \(A_n\) with that of \(C_n\)
in every weight, without a derived-tensor assumption.
Thus \(A\) has the vertical bounds (136.VP1).
The rational common-target theorem ARCH-136-TW-RATIONAL
proves that \(A_{\mathbf Q}\) is admissible, with explicit
uniform horizontal bounds, again without flatness of \(C\).
Apply ARCH-136-VERTICAL-PERIODS.
\hfill\(\square\)

\paragraph{Which tests this actually controls.}
The conclusion applies to every integral closed relative form
already in \(A\), including closed expressions formed using
its multiplication and divided powers. It also controls
products of existing cohomology classes through the
multiplicativity just proved. This is not an assertion of
preservation of specified PD secondary operations by the
associative convolution maps: the latter have no compatible
strict PD structure.

Most importantly, the theorem does not show that adjoining
a chosen rational primitive keeps all the new closed
expressions integral. The explicit zero-form example
ARCH-136-PERIOD-REPAIR-FAILURES shows why testing the
primitive's own endpoint values is insufficient: new
products can have nonintegral periods. An integral lattice
inside a larger rational common target would have to
retain the appropriate tests after every such adjoining
operation. Neither (136.VP2) nor (136.VP3) supplies that
closure or proves that the choices are impossible.

\paragraph{Verification status.}
This is an explicit extraction of the one-sided hypothesis
in the already proved associative comparison, with the
relative extension and nonflat standard-factor consequence
spelled out. The entire saved proof was read by the author
and independently by \texttt{/root/cote1348}, with PASS
in the exact stated scope (GPT-6 Astra Pro; effort not
exposed; 2026-11-22). Root's review, if performed, is
recorded separately.
